\section{Definición del problema}
\label{sec:problema}

\subsection{El riesgo central}

El riesgo de una billetera digital es que el dinero deje de cuadrar. Hay dos
formas de que ocurra:

\begin{description}[leftmargin=2.6cm,style=nextline]
  \item[Pérdida de dinero] Se debita al emisor y, por una falla, nunca se
    acredita al receptor.
  \item[Creación de dinero] Se acredita al receptor dos veces, o se acredita
    sin debitar.
\end{description}

En un sistema real, cualquiera de las dos es un incidente contable. En un
prototipo, es la diferencia entre un proyecto demostrable y una pantalla
bonita sin nada detrás.

\subsection{Los cuatro escenarios que hay que resolver}

\paragraph{Fallo a mitad de la operación.} El proceso muere después de
descontar el saldo del emisor y antes de acreditar al receptor. Si ambas
escrituras no forman una unidad atómica, el dinero desaparece.

\paragraph{Operaciones simultáneas.} Dos peticiones leen el mismo saldo antes
de que ninguna escriba, ambas concluyen que hay fondos suficientes y ambas
debitan. El resultado es un saldo negativo o dinero creado de la nada. Un
ejemplo numérico aparece en el cuadro~\ref{tab:lostupdate}.

\begin{table}[htbp]
\centering
\caption{Actualización perdida con cuentas ficticias. Ana tiene S/\,100 y
recibe dos peticiones simultáneas de S/\,80.}
\label{tab:lostupdate}
\small
\begin{tabular}{@{}clll@{}}
\toprule
$t$ & Petición 1 (a Luis) & Petición 2 (a Rosa) & Saldo en la base \\
\midrule
$t_0$ & lee saldo $=10\,000$ ctms & & $10\,000$ \\
$t_1$ & & lee saldo $=10\,000$ ctms & $10\,000$ \\
$t_2$ & valida $10\,000 \geq 8\,000$ & & $10\,000$ \\
$t_3$ & & valida $10\,000 \geq 8\,000$ & $10\,000$ \\
$t_4$ & escribe $2\,000$ & & $2\,000$ \\
$t_5$ & & escribe $2\,000$ & $2\,000$ \\
\bottomrule
\end{tabular}

\vspace{0.4em}
{\footnotesize Ana envió S/\,160 teniendo S/\,100 y su saldo final es S/\,20.
Se crearon S/\,60 de la nada.}
\end{table}

\paragraph{Respuesta perdida.} El servidor confirma la transacción, pero la
respuesta no llega al cliente por un corte de red. El cliente no sabe si el
dinero se movió. Si reintenta y el servidor trata la petición como nueva, se
paga dos veces. Ninguna propiedad transaccional resuelve esto, porque el
problema ocurre fuera del límite de la transacción.

\paragraph{Acceso a cuentas ajenas.} Si la cuenta de origen se toma del cuerpo
de la petición, cualquiera puede transferir desde la cuenta de otro. Es el
primer riesgo del OWASP API Security Top~10 de 2023~\cite{owaspapi2023}.

\subsection{Un caso real}

El 1 de agosto de 2012, Knight Capital perdió más de 400 millones de dólares
en 45 minutos por un despliegue defectuoso. La orden administrativa de la
Securities and Exchange Commission documenta que la firma no disponía de
salvaguardas adecuadas para limitar los riesgos de su acceso al mercado ni
revisaba la efectividad de sus controles, y se saldó con una sanción de 12
millones de dólares~\cite{sec2013knight}.

El aprendizaje transferible no es sobre bases de datos sino sobre controles:
las invariantes deben ser verificables de forma independiente al código que
las produce. Esa idea es la que justifica que PayCore incluya una prueba de
reconciliación que recalcula los saldos desde el libro mayor y los compara con
la proyección materializada.

\subsection{Alcance}

\textbf{Dentro:} registro y autenticación de usuarios ficticios, consulta de
cuenta y saldo, transferencias entre cuentas, historial de movimientos,
rechazo por fondos insuficientes, prevención de duplicados, procesamiento
correcto de operaciones concurrentes e interfaz web.

\textbf{Fuera, de forma explícita:} dinero real, conexión con entidades
financieras o pasarelas de pago, retiros, recargas desde medios externos,
múltiples divisas, intereses y comisiones.
